% Part: first-order-logic
% Chapter: sequent-calculus
% Section: quantifier-rules

\documentclass[../../../include/open-logic-section]{subfiles}

\begin{document}

\olfileid{fol}{seq}{qrl}

\olsection{量词规则}

\subsection{$\lforall$ 的规则}

\begin{defish}
\Axiom$ !A(t), \Gamma \fCenter \Delta$
\RightLabel{\LeftR{\lforall}}
\UnaryInf$ \lforall[x][!A(x)],\Gamma \fCenter \Delta$
\DisplayProof
\hfill
\Axiom$ \Gamma \fCenter \Delta, !A(a) $
\RightLabel{\RightR{\lforall}}
\UnaryInf$ \Gamma \fCenter \Delta, \lforall[x][!A(x)]$
\DisplayProof
\end{defish}

在 \LeftR{\lforall} 中，$t$ 是一个闭项（即不含变元的项）。在 \RightR{\lforall} 中，$a$~是一个常项符号，且不得出现在 \RightR{\lforall} 规则下相继式中的任何位置。我们称 $a$ 为 \RightR{\forall} 推理的\emph{特征变元}。\footnote{尽管上述规则中的 $a$ 是常项符号，我们仍沿用“特征变元”这一术语，这是出于历史原因。}

\subsection{$\lexists$ 的规则}

\begin{defish}
\Axiom$ !A(a), \Gamma \fCenter \Delta $
\RightLabel{\LeftR{\lexists}}
\UnaryInf$ \lexists[x][!A(x)], \Gamma \fCenter \Delta$
\DisplayProof
\hfill
\Axiom$ \Gamma \fCenter \Delta, !A(t) $
\RightLabel{\RightR{\lexists}}
\UnaryInf$ \Gamma \fCenter \Delta, \lexists[x][!A(x)]$
\DisplayProof
\end{defish}

同样，$t$~是闭项，$a$~是常项符号，且不得出现在 \LeftR{\lexists} 规则的下相继式中。我们称 $a$ 为该 \LeftR{\lexists} 推理的\emph{特征变元}。

特征变元不在 \RightR{\lforall} 或 \LeftR{\lexists} 推理的下相继式中出现的条件，称为\emph{特征变元条件}。

\begin{explain}
回忆一下我们的记法约定：当公式 $!A$ 含有自由变元~$x$ 时，我们记作 $!A(x)$。在此上下文中，$!A(t)$ 则是 $\Subst{!A}{t}{x}$ 的简写。因此，我们也可以将 \RightR{\lexists} 规则写作：
\begin{prooftree}
  \Axiom$\Gamma \fCenter \Delta, \Subst{!A}{t}{x}$
  \RightLabel{\RightR{\lexists}}
  \UnaryInf$\Gamma \fCenter \Delta, \lexists[x][!A]$
\end{prooftree}
请注意，$t$ 可能已在 $!A$ 中出现，例如 $!A$ 可能是 $\Atom{\Obj P}{t,x}$。因此，从 $\Gamma \Sequent \Delta,
\Atom{\Obj P}{t,t}$ 推出 $\Gamma \Sequent \Delta,
\lexists[x][\Atom{\Obj P}{t,x}]$ 是 \RightR{\lexists} 的一个正确应用——你可以用约束变元 $x$ “替换”前提中一个或多个（不必全部）$t$ 的出现。然而，\RightR{\lforall} 和 \LeftR{\lexists} 中的特征变元条件要求常项符号 $a$ 不在 $!A$ 中出现。所以，你不能正确地使用 $\RightR{\lforall}$ 从 $\Gamma \Sequent \Delta, \Atom{\Obj P}{a,a}$ 推出 $\Gamma \Sequent \Delta, \lforall[x][\Atom{\Obj P}{a,x}]$。
\end{explain}

\begin{explain}
在 \RightR{\lexists} 和 \LeftR{\lforall} 中，对项 $t$ 没有限制。另一方面，在 \LeftR{\lexists} 和 \RightR{\lforall} 规则中，特征变元条件要求常项符号 $a$ 不能在上相继式中 $!A(a)$ 之外的任何位置出现。这一条件是确保系统可靠所必需的，即系统只推导有效的相继式。如果没有这个条件，下面这样的推导就会被允许：
\begin{prooftree}
  \Axiom$!A(a) \fCenter !A(a)$
  \RightLabel{*\LeftR{\lexists}}
  \UnaryInf$\lexists[x][!A(x)] \fCenter !A(a)$
  \RightLabel{\RightR{\lforall}}
  \UnaryInf$\lexists[x][!A(x)] \fCenter \lforall[x][!A(x)]$
  \DisplayProof\bottomAlignProof
  \qquad
  \Axiom$!A(a) \fCenter !A(a)$
  \RightLabel{*\RightR{\lforall}}
  \UnaryInf$!A(a) \fCenter \lforall[x][!A(x)]$
  \RightLabel{\LeftR{\lexists}}
  \UnaryInf$\lexists[x][!A(x)] \fCenter \lforall[x][!A(x)]$
\end{prooftree}
然而，$\lexists[x][!A(x)] \Sequent \lforall[x][!A(x)]$ 是无效的。
\end{explain}

\end{document}